\section{Limitations and next steps}
The implementation has the following scope limits:
\begin{itemize}
\item Native and Torch paths have different operator, dtype and gradient coverage.
Traced modules support a guarded subset; structured reverse loops and general
higher-order differentiation remain incomplete.
\item Generated exact kNN supports a bounded FP32, small-k envelope. Larger-k
scratch, general metric lowering and generated backward require further work.
\item Implicit block composition, automatic heterogeneous out-of-core execution
and additional vendor backends remain incomplete.
\item Distributed execution uses fixed ownership. Compute-aware partitioning,
partition-local topology capacity, RCCL and failure recovery remain open.
\item Gaussian and volume rendering use host geometry work and are not differentiable.
\item Evaluation measures FP32 forward workloads. Sustained out-of-core training
and beyond-host-RAM graphs require additional measurement.
\end{itemize}

\section{Conclusion}
Tiga retains relations, reducer algebra, gradients and execution constraints as
a compiler-level contract. Generated relations can avoid costly materialization;
this improves radius and exact kNN latency and allocation on the measured
workloads. Paging permits a fully checked billion-edge
forward evaluation on one 16 GiB GPU, while profiling exposes substantial host
staging cost. Two-host execution works but does not beat the faster single GPU
under the measured equal-row schedule. The resulting priorities are concrete:
reduce dynamic dispatch cost, reduce staging copies, and make distributed
scheduling sensitive to compute imbalance and halo traffic.
